\chapter{LITERATURE REVIEW}
\label{ch:litreview}
% Promoted out of Chapter 1 on 20 Sep 2026: thesis.md lists the literature
% review as a chapter in its own right, with four strands it names explicitly.
% Sources it also names:
% 20 Sep 2026, author's instruction: ALL complexity-theory and Turing-machine
% material was DROPPED from this chapter -- the former Sec. 2.1 "Models of
% Computation and the Limits of Classical Computing", Sec. 2.2 "Quantum
% Computing and Computational Complexity" and Sec. 2.3 "Quantum Walks as a
% Physical Realisation of Non-Deterministic Computation" were deleted whole,
% together with the chapter's complexity-class table and containment chain.
% Do not mine thesis.tex for automata/Turing prose, and do not copy the
% complexity-class comparison across from outdated__only_for_reference.tex:
% both were sources for exactly the material that has been removed.
% Prose written 20 Sep 2026. Two standing rules for this chapter:
%   * NO ITALICS anywhere in the running text (no \emph, no \textit). Where
%     prose was carried across from thesis.tex or from
%     outdated__only_for_reference.tex, the \emph calls in the source were
%     removed rather than reproduced.
%   * QUANTUM REFERENCES ONLY. The former Sec. 2.5, "Machine Learning:
%     Classical Foundations", was deleted on author's instruction: cite no
%     purely classical machine-learning work here. The classical treatment
%     belongs to Ch. 3 Sec. "An Overview of Machine Learning", which the spec
%     already provides for and which the appendix listings support. Removed
%     with it: Samuel_1959, Awad_2015, alpaydin2021machine, learning1997tom,
%     ElNaqa2015, Hofmann_2008, Valkenborg_2023, Xia_2020, Jackson_2014,
%     Perozzi_2014, Grover_2016, SotoGomez_2021, Bozorgi_2024, Xu_2021,
%     doi:10.1137/20M1386062, Krzakala_2013, Vaswani_2017, Sen_2008,
%     Macskassy_2007, Bttcher_2022. All remain in the .bib.

This chapter surveys the literature on which the thesis rests, organised as a
single argument rather than as a catalogue.
Section~\ref{sec:lr-primitive} reviews the quantum-walk literature proper and
Section~\ref{sec:lr-qml} the quantum machine-learning literature, with
Section~\ref{sec:lr-dtqwqml} covering the small body of work that combines the
two. Section~\ref{sec:lr-gap} states the gap the thesis addresses. Classical
machine learning is treated in Chapter~\ref{ch:prelim} and is not reviewed
here; the references in this chapter are to quantum work throughout.


\section{Quantum Walks as an Algorithmic Primitive}
\label{sec:lr-primitive}
% thesis.md strand 4 begins here: the walk literature proper.

\subsection{Discrete- and Continuous-Time Walks in the Literature}
The quantum walk was introduced in two independent forms. Aharonov, Davidovich
and Zagury~\cite{Aharonov_1993} formulated the discrete-time walk as the
repeated application of a coin operator on an internal degree of freedom
followed by a shift conditioned on it; Farhi and Gutmann~\cite{Farhi_1997}
formulated the continuous-time walk as evolution under a Hamiltonian given by
the adjacency matrix of a graph, with no coin register at all. The two were
developed separately for a decade and later shown equivalent as models of
computation~\cite{Childs2010, strauch2006connecting}, with hybrid formulations
unifying them proposed more recently~\cite{chen2025hybrid}.

The behaviour distinguishing the quantum walk from its classical counterpart was
established early. On the line the quantum walker spreads ballistically, its
standard deviation growing linearly in the number of steps rather than as the
square root, and its limiting distribution is bimodal rather than
Gaussian~\cite{Ambainis_2001, Konno_2002, Grimmett_2003}; Knight, Roldan and
Sipe~\cite{Knight_2003} showed this to be an interference phenomenon requiring
no quantisation of the walker as such. On graphs, Aharonov et
al.~\cite{Aharonov_2000} established the general framework and bounded the
mixing time, and Childs et al.~\cite{Childs_2001, Childs_2002} exhibited a
glued-trees graph on which the continuous-time walk traverses from one root to
the other exponentially faster than any classical algorithm, which remains the
strongest separation known for a walk. Watrous~\cite{Watrous_2001} connected
walks to space-bounded simulation of classical walks, and comprehensive reviews
are available~\cite{Elias_2012, Kadian_2021, Xiaogang_2024, Mlken_2011}.

The coin is the feature of the discrete-time walk this thesis is concerned with,
and the literature on it is substantial but has a consistent shape. Tregenna et
al.~\cite{Tregenna_2003} showed that the coin and the initial state jointly
control the shape of the distribution; Brun, Carteret and
Ambainis~\cite{Brun_2002} studied walks driven by many coins; Patel et
al.~\cite{Patel_2004} showed the coin can be dispensed with entirely. Beyond the
fixed real coins that are standard --- Hadamard, Grover and
Fourier~\cite{Komatsu_2019, Saito_2018} --- coins have been made dependent on
time~\cite{C_2006}, on the step index~\cite{Panahiyan_2018}, on
position~\cite{Ahmad_2020} and on edge weights~\cite{Wong_2017}; general $U(2)$
coins have been parametrised and swept~\cite{Li_2012, Souza_2013, Naves_2023};
inhomogeneous coins produce Parrondo-type
effects~\cite{Mittal_2024, Kadiri_2024}; and coin sequences have been optimised
for entanglement generation~\cite{Gratsea_2020}. This body of work establishes
that the coin materially changes the walk. It does not connect that freedom to a
learning objective, which is the gap Section~\ref{sec:lr-gap} states.

Alternative formulations recur later. Szegedy's quantisation~\cite{Szegedy}
builds a walk from an arbitrary reversible Markov chain and provides the
standard route to quadratic speed-ups over classical mixing, and is equivalent
to the coined walk under mild conditions~\cite{G_2017}. The staggered
model~\cite{Portugal_2016} dispenses with the coin in favour of graph
tessellations, and the lackadaisical walk~\cite{G_2018} adds a weighted
self-loop at each vertex.

\subsection{From Search to Universality}
Search is the application in which walk-based algorithms are best understood.
Shenvi, Kempe and Whaley~\cite{Shenvi_2002} gave a discrete-time walk on the
hypercube finding a marked vertex in $O(\sqrt{N})$ steps, matching Grover;
Childs and Goldstone~\cite{Childs_2003, Childs_2004} analysed spatial search by
continuous-time walk and established the dimensional thresholds at which the
speed-up survives; Ambainis, Kempe and Rivosh and later
Tulsi~\cite{Ambainis_2004, Tulsi_2008} improved the two-dimensional case;
Magniez et al.~\cite{Magniez_2007} gave the general framework; and surveys are
available~\cite{santha2008quantum}.

Subsequent work has broadened the setting considerably. Search has been analysed
for multiple marked vertices~\cite{A_2021, bezerra2021quantum} and made robust
to ignorance of their number~\cite{Xu_2021a}; deterministic variants using
alternating walks have been constructed~\cite{Marsh_2021}; the quadratic
speed-up in finding, as opposed to detecting, marked vertices has been
established in generality~\cite{Ambainis_2019} and matched in continuous
time~\cite{Apers_2021}; and the effect of multiple walk steps per oracle query
has been characterised~\cite{wong2015quantum}. A distinct thread concerns the
lackadaisical walk, in which a weighted self-loop lets the walker remain in
place~\cite{G_2018}, improving the success probability of search on
vertex-transitive~\cite{Rhodes_2020} and complete bipartite~\cite{Rhodes_2018}
graphs; the exceptional configurations in which walk-based search fails outright
have been catalogued~\cite{Wong_2017a}.

The sub-literature around the weight of that self-loop deserves separate
mention, because it is the clearest demonstration in the search literature that
a single scalar parameter of the walk operator governs whether a speed-up is
obtained at all. Giri and Korepin~\cite{Giri_2020} found that only specific
weights produce the speed-up on one- and two-dimensional lattices;
Nahimovs~\cite{Nahimovs_2018} showed that the weight established for a single
marked vertex on the two-dimensional grid ceases to be optimal once several are
marked; Rhodes and Wong~\cite{Rhodes_2019} dropped the assumption that every
self-loop carries the same weight, which is natural only on vertex-transitive
graphs; and Rapoza and Wong~\cite{Rapoza_2021} showed that the self-loops may
break every symmetry of a vertex-transitive graph without loss of speed-up.

The pattern recurs in the learning applications reviewed later: the performance
of a quantum-walk algorithm is frequently decided not by the graph, nor by the
broad architecture of the walk, but by a small number of parameters internal to
the coin, whose correct values are instance-dependent. More recent proposals
apply Grover search within the coin space itself~\cite{Giri_2023}, construct
search-complement algorithms~\cite{WingBocanegra_2025}, guide the walk towards
the target~\cite{Schulz_2024} or place the marked vertex in
motion~\cite{Sahu_2024}, and the connection to vertex centrality was noted early
by Berry and Wang~\cite{Berry_2010}. A constant-time algorithm for unstructured
search using a many-body continuous-time walk has been
proposed~\cite{DalFavero_2024}; it is a continuous-time result with no
established discrete-time analogue, and Section~\ref{sec:search} records it
as such; Section~\ref{sec:res-search} quantifies the trade-off.

\subsection{Walks Beyond Search: Optimisation and Algebraic Problems}
The quantum walk is not confined to search. Ambainis~\cite{Ambainis_2003} gave a
walk on the Johnson graph solving element distinctness in $O(N^{2/3})$ queries,
matching the known lower bound and resolving a problem that had resisted
amplitude-amplification approaches; the technique generalises to problems in
which a costly checking step is amortised against many cheap update
steps~\cite{Ambainis_2004}. Childs et al.~\cite{Childs_2007} applied a
discrete-query walk to the evaluation of NAND formulae, and walk-based
techniques have since been applied across algebraic
problems~\cite{Childs_2008} --- the connection Section~\ref{sec:shor}
follows in setting out what a walk-based reformulation of Shor's algorithm
does and does not amount to.

In optimisation, the quantum approximate optimisation algorithm
(QAOA)~\cite{Farhi_2014, Farhi_2016} alternates a problem Hamiltonian with a
mixing Hamiltonian, and the mixer is a continuous-time walk on the hypercube;
that identification is what makes QAOA a walk algorithm rather than merely an
algorithm resembling one, and the structure is developed in
Section~\ref{sec:optimisation}. Tutorial and review treatments are
available~\cite{Jaeho_2019, Chicano_2025}, and the symmetry and Lie-algebraic
structure of the ansatz has been analysed~\cite{Kazi_2025}. Explicitly
walk-based optimisation has been developed as the quantum walk optimisation
algorithm~\cite{Marsh_2020}, benchmarked against classical heuristics on
weighted maximum cut~\cite{Bennett_2025}, formulated in discrete
time~\cite{Liliopoulos_2024}, and applied to sampling through a quantum
Metropolis construction~\cite{Campos_2023}. Quantum random access memory has
also been formulated as a walk on a binary
tree~\cite{Asaka_2021, Asaka_2023, Mullor_2022}, an application relevant to the
data-loading bottleneck of Section~\ref{sec:lr-qml}.

The application areas this thesis keeps in view are the higher-level ones. Walks
underlie proposals for entanglement distribution across arbitrary quantum
networks~\cite{Chen_2024c} and for multi-node quantum
teleportation~\cite{Ikken_2025}, both quantum-internet primitives; complex
quantum networks are reviewed by Nokkala et al.~\cite{Nokkala_2024} and the
communication-capacity consequences of superposed paths by Caleffi et
al.~\cite{Caleffi_2023}. In routing and logistics, a path-integral formulation
closely related to the walk has been applied to vehicle routing under a qubit
budget~\cite{Gautam_2024}; in cryptography, walks supply certified random-number
generation~\cite{Haider_2023}. These are the use cases the synopsis-defence
committee asked to see connected to the probabilistic-machine framing, and they
are why the walk is treated here as a general computational tool rather than a
graph-search heuristic.

Circuit implementation matters for the numerical work reported later. Efficient
constructions for coined walks were given by Douglas and
Wang~\cite{Douglas_2007} and refined subsequently~\cite{Han_2017}, with scalable
implementations based on the quantum Fourier transform proposed by
Shakeel~\cite{Shakeel_2020}; constructions have been extended from regular
graphs to complex networks~\cite{Sato_2024, Sato_2025} and, in the
continuous-time case, to sparse~\cite{Chen_2024a} and
random~\cite{Chakraborty_2025} graphs, with multi-qubit computation on closed
graphs set out by Chawla et al.~\cite{Chawla_2023}. The hardware constraints
under which any of this runs are given by Preskill~\cite{Preskill_2018} and by
materials-level surveys~\cite{doi:10.1126/science.abb2823}.


\section{Quantum Machine Learning}
\label{sec:lr-qml}

\subsection{Data Encoding and Feature Maps}
Quantum machine learning is the study of learning algorithms that employ quantum
processing in some part of their operation. Introductions and surveys spanning
the field's development are
available~\cite{Schuld_2015, Biamonte_2017, Ciliberto_2018, Tychola_2023,
Zeguendry_2023}, with a critical assessment by Cerezo et
al.~\cite{Cerezo_2023} and a course-level treatment by IBM
Quantum~\cite{IBMQML2026}.

Early enthusiasm rested substantially on the algorithm of Harrow et
al.~\cite{Harrow_2008} for linear systems, which appeared to offer an
exponential speed-up for a primitive underlying much of learning theory. Its
qualifications --- that the matrix be sparse and well-conditioned, that the
input be available in quantum superposition, and that only a summary statistic
of the solution be required --- proved more restrictive in learning applications
than was initially appreciated. The general difficulty is data access: an
algorithm whose quantum subroutine runs in time polylogarithmic in the input
size confers no advantage if loading the input costs time linear in it. This
motivates the walk-based quantum random access memory proposals noted
earlier~\cite{Asaka_2021}, and it is why the field has moved towards algorithms
that consume classical data through a narrow interface.

That interface is the feature map. Encoding a classical vector into a quantum
state is a feature map, and the choice of encoding fixes the function class the
model can represent: Schuld, Sweke and Meyer~\cite{Schuld_2021a} showed that a
variational model with repeated data encoding is a truncated Fourier series
whose accessible frequency spectrum is determined by the encoding Hamiltonian,
not by the trainable parameters. This is why the encoding rather than the ansatz
is the first design decision in Chapter~\ref{ch:methodology}, and it is what
makes a walk a candidate feature map at all: the walk operator is an encoding
with graph structure built into it.

\subsection{Quantum Neural Networks}
\label{sec:lr-qnn}
Under the designation quantum neural network, parametrised circuits have been
studied as learning models in their own right~\cite{Beer_2022, Zhou_2023}. Their
capacity has been analysed through the Fisher information spectrum, on which
basis Abbas et al.~\cite{Abbas_2021} argued that well-designed quantum models
can exhibit higher effective dimension and faster training than comparable
classical networks. Architectures proposed include data re-uploading schemes in
which a single qubit becomes a universal classifier by interleaving data
encoding with trainable rotations~\cite{PrezSalinas_2019}, variational state
discriminators~\cite{Lee_2023}, generative models~\cite{Gao_2018}, equivariant
networks whose symmetry is matched to that of the data~\cite{Nguyen_2024} and
fully quantum optimisation subroutines~\cite{Chen_2024b}. Architecture search
has itself been approached with reinforcement learning~\cite{Rapp_2024}, and
experimental evaluations on present hardware reported~\cite{Simoes_2023}.

The principal obstacle is trainability. McClean et al.~\cite{McClean_2018}
showed that for a broad class of random parametrised circuits the gradient of
the cost function vanishes exponentially in the number of qubits, so the
optimisation landscape is a barren plateau on which gradient-based training
fails. This is not peripheral: it constrains circuit depth, ansatz choice and
cost-function locality simultaneously, and any proposal for a new quantum model
--- including a walk-based one --- must be assessed against it.

\subsection{Quantum Variational Circuits}
The dominant paradigm in the noisy intermediate-scale era~\cite{Preskill_2018}
is the variational quantum circuit: a parametrised circuit whose gate parameters
are optimised by a classical routine against a measured cost function. The
approach was set out as quantum circuit learning by Mitarai et
al.~\cite{Mitarai_2018}, framed as a model class by Benedetti et
al.~\cite{Benedetti_2019}, and is reviewed comprehensively by Cerezo et
al.~\cite{Cerezo_2021}. The designations quantum neural network and quantum
variational circuit are used almost interchangeably in the literature, and are
treated as one family here, differing in emphasis rather than substance.

Because the ansatz is a free choice, comparing candidate circuits requires
structural measures. Sim, Johnson and Aspuru-Guzik~\cite{Sim_2019} introduced
expressibility, the extent to which a circuit's output states cover the Hilbert
space, together with entangling capability, and showed that the two trade off
against trainability. The relevance here is direct: a walk step is a particular
choice of ansatz with a fixed entangling pattern given by the graph, and it must
be justified against these measures rather than assumed superior because the
walk has good properties elsewhere.

\subsection{Quantum Kernel Methods}
\label{sec:lr-qkernel}
A learning algorithm that depends on its data only through inner products can be
applied in a high-dimensional feature space without constructing the feature map
explicitly, provided the inner product there can be evaluated as a kernel on the
original data. That structure gives exactly the narrow data interface described
above, and makes kernels the most theoretically satisfactory part of quantum
machine learning. Schuld and Killoran~\cite{Schuld_2018} observed that encoding
a classical data point into a quantum state is a feature map, that the inner
product of two such states is therefore a kernel, and that a quantum computer
may consequently serve as a kernel evaluator. Havlicek et
al.~\cite{Havlicek_2019} implemented both a variational classifier and a kernel
estimator on superconducting hardware in one work, which is why that paper is
the common ancestor of both families. Overlaps of the form required are
estimated by the swap test~\cite{Buhrman_2001} or by the inversion circuit, and
covariant kernels adapted to data with group structure were developed
subsequently~\cite{Glick_2021}. Schuld~\cite{Schuld_2021} later argued that
supervised variational models are themselves kernel methods, unifying
Sections~\ref{sec:lr-qnn} to~\ref{sec:lr-qkernel} under a single analysis.

This line of work has produced the field's clearest positive and negative
results, and both should be stated. Positively, Liu et al.~\cite{Liu_2020}
constructed a classification problem, based on the discrete logarithm, for which
a quantum kernel support vector machine achieves a rigorous speed-up requiring
only classical access to the data. Negatively, Huang et al.~\cite{Huang_2021}
showed that the availability of training data substantially narrows the class of
problems on which an advantage can persist; Thanasilp et
al.~\cite{Thanasilp_2024} showed that quantum kernel values can concentrate
exponentially in the number of qubits, so that distinguishing them requires
exponentially many shots and the apparent advantage evaporates; and the
complexity of training quantum support vector machines has been
analysed~\cite{Gentinetta_2022}. Quantum kernels are thus well understood and
the conditions for advantage correspondingly narrow. Exponential concentration
is in particular a design constraint on any walk-based kernel, since a deep walk
is exactly the kind of expressive encoding that induces it.

\subsection{Benchmarking and the Question of Quantum Advantage}
\label{sec:lr-bench}
% thesis.md points at PennyLane's qml-benchmarks collection and IBM Quantum's
% QML course as implementation references; the benchmarking literature is the
% honest frame for any comparison this thesis makes against classical models.
The evaluation practices of the field have come under sustained criticism, and
the criticism is accepted here. Schuld and Killoran~\cite{Schuld_2022} argued
that quantum advantage is the wrong organising goal for quantum machine
learning, since the comparison is almost always against an under-tuned baseline
on a dataset chosen after the fact. Bowles, Ahmed and Schuld~\cite{Bowles_2024}
made the point quantitatively: across a systematic benchmark of twelve quantum
models on a curated dataset collection, no quantum model outperformed a
well-tuned classical baseline, and apparent advantages traced to hyperparameter
selection and dataset choice. Their dataset collection is the benchmark suite
this thesis uses, and their methodological requirements --- a tuned baseline, a
fixed evaluation split chosen before the model is selected, and results reported
across seeds --- are adopted in Chapter~\ref{ch:results}. Cerezo et
al.~\cite{Cerezo2022} survey the challenges and opportunities on the same terms.

The consequence for this thesis is stated plainly. The simulations reported
later run on classical simulators and establish relative behaviour between walk
configurations, not advantage over classical machine learning. No claim of
quantum advantage is made anywhere in this work.


\section{Discrete-Time Quantum Walks in Quantum Machine Learning}
\label{sec:lr-dtqwqml}

\subsection{Walk-Based Neural and Variational Models}
The literature combining the two subjects is considerably smaller than either.
Schuld, Sinayskiy and Petruccione~\cite{Schuld_2014} gave the earliest proposal,
representing the firing patterns of a quantum neural network as a walk on a
graph. Dernbach et al.~\cite{Dernbach_2018, Dernbach_2019} introduced quantum
walk neural networks for graph-structured data and, in the later work, made the
coin depend on learned node features, which is the closest antecedent of the
present thesis and is discussed again in Section~\ref{sec:lr-gap}. de Souza et
al.~\cite{deSouza_2022} used a quantum walk as the search procedure for training
a neural network, inverting the usual arrangement, and Zhu et
al.~\cite{Zhu_2026} recently showed that a quantum graph neural network can be
realised on a single qubit using a walk. Casares, Campos and
Martin-Delgado~\cite{Casares_2022} combined walks with deep learning for protein
folding, one of the few end-to-end applications in the area. Critical reviews of
quantum graph neural networks~\cite{Ceschini_2026} and of how embeddings shape
them~\cite{innan2026embeddingsshapegraphneural} have appeared very recently, and
both conclude that the empirical case is not yet made.

\subsection{Walk-Based Kernels}
Graph kernels built from walks predate the variational literature. Rossi,
Torsello and Hancock~\cite{Rossi_2013} constructed a continuous-time walk kernel
for unattributed graphs, and Zhang et al.~\cite{Zhang_2022} gave an
R-convolution kernel based on a fast discrete-time walk. These are kernels
between graphs, computed from walk statistics, and should be distinguished from
the quantum kernels of Section~\ref{sec:lr-qkernel}, which are inner products
between quantum states evaluated on quantum hardware. The combination --- a walk
used as the state-preparation circuit of a quantum kernel, so that the kernel
entry is the overlap of two walk-evolved states --- is the construction
Chapter~\ref{ch:methodology} develops, and the literature on it is thin.

\subsection{Walk-Based Representation Learning}
% Where node embedding and QWalkVec belong: related work only. Node
% embedding was dropped from the thesis on 20 Sep 2026.
Node embedding by quantum walk has attracted the most attention of the three.
QWalkVec~\cite{10.1007/978-981-97-2242-6_8} substitutes a coined quantum walk
for the classical walk used by standard node-embedding methods and reports
improved downstream classification. RED~\cite{Wang_2021} learns role embeddings
by discrete-time walk, and GraphQWalk~\cite{Liu_2026} the structural analogue by
continuous-time walk. Walks have also been applied to community
detection~\cite{Mukai_2020, He_2025}, to nearest-neighbour
classification~\cite{Enrico_2024}, to centrality and critical-node
identification~\cite{Wang_2025, Liang_2022, Wang_2022, Berry_2010} and to the
degree structure of networks~\cite{Faccin_2013}, and the scaling of walk search
on complex networks has been characterised~\cite{Sato_2024b}.

Two qualifications apply to this sub-literature and are carried into
Chapter~\ref{ch:results}. First, the reported gains are typically small,
obtained on a handful of small graphs, and sensitive to the truncation length
and to the downstream classifier, neither of which is always held fixed across
the compared methods. Second, walks with memory --- a register recording some
part of the walker's history~\cite{Flitney_2003, McGettrick_2009, Gettrick_2014,
Krawec_2014, Zhou_2019, Li_2020a, Shakeel_2019, Camilleri_2014,
PhysRevA.87.052302, PhysRevA.106.062215, 2011451.2011461} --- have been proposed
as the quantum counterpart of the second-order classical walk. They are recorded
here as related work only; they are not part of the contribution of this thesis.


\section{Quantum Walks as a Comprehensive Algorithmic Framework for Quantum Machine Learning: the Gap Addressed by This Thesis}
\label{sec:lr-gap}
% The gap argument, restated for the narrowed contribution: walk-based
% learning algorithms fix the coin by convention -- Hadamard, Grover or
% Fourier -- and treat that choice as given. The coin is in fact a design
% space: its weights may depend on local graph structure and on the step
% index, and its entries need not be real. What that freedom is worth to a
% walk-based learner has not been established. Memory-based walks are not
% part of this argument; mention them as related work or not at all.

Three observations follow from the survey above, and together they define the
gap this thesis addresses.

First, the walk literature of Section~\ref{sec:lr-primitive} establishes
repeatedly that the performance of a walk-based algorithm is governed by a small
number of parameters internal to the coin, and that the correct value of those
parameters is instance-dependent. The lackadaisical self-loop weight is the
clearest case: a scalar that decides whether a quadratic speed-up is obtained at
all, whose optimum shifts with the number of marked
vertices~\cite{Nahimovs_2018}, need not be uniform across
vertices~\cite{Rhodes_2019}, and may break the symmetry of the graph
entirely~\cite{Rapoza_2021}. The coin is thus known within the walk literature
to be a design space rather than a convention.

Second, the walk-based learning literature of Section~\ref{sec:lr-dtqwqml} does
not treat it as one. With the partial exception of Dernbach et
al.~\cite{Dernbach_2019}, whose coin depends on learned node features,
walk-based learning models fix the coin in advance to one of Hadamard, Grover or
Fourier and report results for that choice. The coin is not made dependent on
local graph structure, nor on the step index, and its entries are taken to be
real. Each of these three freedoms is standard in the walk literature ---
position-dependent~\cite{Ahmad_2020}, step-dependent~\cite{Panahiyan_2018} and
general $U(2)$~\cite{Li_2012, Souza_2013} coins are all well studied --- and
none has been evaluated against a learning objective.

Third, the benchmarking literature of Section~\ref{sec:lr-bench} shows why the
gap has gone unnoticed. Where a walk-based model is compared only against a
weak baseline and reports a win, there is no reason to ask whether a different
coin would have done better; and where the comparison is against a tuned
baseline, the walk-based models have not so far won~\cite{Bowles_2024}. The coin
has thus been neither optimised nor ruled out.

The thesis accordingly asks a narrower question than the chapter title might
suggest. Taking the discrete-time coined walk as a computational primitive, and
instantiating it in the three mainstream quantum machine learning families of
Section~\ref{sec:lr-qml} --- quantum neural networks, quantum variational
circuits and quantum kernel methods --- what does treating the coin as a design
space buy? Two axes are varied: whether the coin depends on local graph
structure and on the step index, and whether its entries are permitted to be
complex. The question is empirical and is answered on simulators, against tuned
classical baselines, under the reporting discipline of Bowles et
al.~\cite{Bowles_2024}. A negative answer on either axis is a result and is
reported as one.


